\begin{textsample}{POS Dim 6 – pb – Score 12.00 – pb/pb\_em\_94.txt}  \label{ex:f6_pos_201}
3.3.2 \textbf{PERFIL} DO ESTUDANTE EGRESSO Dos estudantes que fizerem a opção pelo Itinerário \textbf{Formativo} TSS, esperamos que, ao término dessa etapa de ensino, sejam capazes de contribuir com soluções para os mais diversos enfrentamentos ; e para o desenvolvimento e/ou o melhoramento de novas \textbf{tecnologias} para o setor de serviços e de bens e consumo, de forma crítica em relação às consequências do emprego das \textbf{tecnologias}, hoje existentes, para o meio ambiente.
Além disso, ter autonomia para empregar os conhecimentos na idealização de projetos pessoais com vistas a melhoria da qualidade de vida individual e coletiva.
Também esperamos que o egresso tenha capacidade de se perceber não só como beneficiário, mas também como criador de produtos, bens e serviços, podendo elaborar projetos e executar atividades fundamentados cientificamente e, sobretudo, com responsabilidade ética, financeira, \textbf{ambiental} e social. 3.3.3 OBJETIVOS GERAIS DO ITINERÁRIO \textbf{FORMATIVO} - Desenvolver as habilidades de criação de serviços, sistemas e \textbf{tecnologias} sustentáveis com base nos Eixos \textbf{Estruturantes} : Investigação Científica, Processos Criativos, \textbf{Empreendedorismo} e \textbf{Mediação} e Intervenção Cultural. - Capacitar o estudante quanto ao aprimoramento da percepção de situações-problema, e na análise de variáveis que interferem na dinâmica de fenômenos da natureza e/ou de processos tecnológicos, considerando dados e informações disponíveis em diferentes mídias, com ou sem o uso de dispositivos e \textbf{aplicativos} \textbf{digitais}. - Desenvolver a habilidade de seleção e sistematização de dados e informações ( independentemente da fonte ), seja no aspecto pessoal quanto no coletivo, promovendo assim a otimização da aprendizagem, quanto ao desenvolvimento do protagonismo e proativismo individual, no sentido da geração de hipóteses e soluções frente aos \textbf{problemas} do cotidiano. - Evidenciar a relevância da investigação científica como instrumento de validação, ou não, de hipóteses, frente às situações-problema ( variáveis ) que interferem na dinâmica de fenômenos da natureza e/ou de processos tecnológicos, com ou sem o uso de dispositivos e \textbf{aplicativos} \textbf{digitais}, utilizando procedimentos e linguagens adequados à investigação científica. - Motivar os estudantes quanto a relevância de métodos alternativos de resolução de conflitos, utilizando para isso a \textbf{mediação} e outras formas de intervenção na resolução de \textbf{problemas} de natureza sociocultural e de natureza \textbf{ambiental} relacionados às Ciências Exatas e da Natureza. - Capacitar o estudante quanto ao uso intencional de conhecimentos e recursos das Ciências Exatas e da Natureza para propor ações individuais e/ou coletivas de \textbf{mediação} e intervenção sobre \textbf{problemas} socioculturais e \textbf{problemas} \textbf{ambientais}. - Desenvolver a capacidade de intervir em produtos, serviços e \textbf{tecnologias} buscando a \textbf{sustentabilidade} e eficiência na utilização de recursos, assim como as características do contexto local e global. - Fornecer subsídios para o exercício da cidadania de forma autônoma, crítica e participativa.

% matched lemmas: ambiental, aplicativo, digital, empreendedorismo, estruturante, formativo, mediação, perfil, problema, sustentabilidade, tecnologia
\end{textsample}
